\documentclass[12pt,a4paper]{article}

% =========================================================
% PACOTES
% =========================================================

\usepackage[utf8]{inputenc}
\usepackage[T1]{fontenc}
\usepackage[brazil]{babel}
\usepackage{tcolorbox}
\usepackage{mathptmx}


\usepackage{geometry}
\geometry{
    a4paper,
    left=2.5cm,
    right=2.5cm,
    top=2.5cm,
    bottom=2.5cm
}

\usepackage{xcolor}
\usepackage{titlesec}
\usepackage{graphicx}
\usepackage{setspace}
\usepackage{hyperref}
\usepackage{fancyhdr}

% =========================================================
% ESTILO TERMINAL / CÓDIGO
% =========================================================

\usepackage{listings}

\definecolor{terminalbg}{RGB}{12,12,12}
\definecolor{terminalgreen}{RGB}{0,255,120}
\definecolor{terminalframe}{RGB}{0,180,90}

\lstdefinestyle{ravenstyle}{
    backgroundcolor=\color{terminalbg},
    basicstyle=\ttfamily\small\color{terminalgreen},
    breaklines=true,
    frame=single,
    rulecolor=\color{terminalframe},
    framesep=8pt,
    tabsize=4,
    showstringspaces=false,
    numbers=left,
    numberstyle=\tiny\color{terminalgreen},
    keywordstyle=\color{green},
    commentstyle=\color{terminalgreen},
    stringstyle=\color{terminalgreen},
    xleftmargin=10pt,
    xrightmargin=10pt,
    keepspaces=true
}

% =========================================================
% CABEÇALHO
% =========================================================

\pagestyle{fancy}

\fancyhf{}

\fancyhead[L]{Projeto Raven / UBR Totem}
\fancyhead[R]{Documentação Técnica}
\fancyfoot[C]{\thepage}

% =========================================================
% TÍTULO
% =========================================================

\title{
    \Huge \textbf{DOCUMENTAÇÃO TÉCNICA}\\[0.4cm]
    \Large Unidade Básica RAVEN (Rede Adaptativa de Variáveis em Estado Não-linear) "TOTEM"
}

\author{
    Vicenzo "Codas" Minossi
}

\date{\today}

% =========================================================
% DOCUMENTO
% =========================================================

\begin{document}

\maketitle

\vspace{1cm}

\tableofcontents

\newpage

% =========================================================
% IDENTIFICAÇÃO
% =========================================================

\section{Identificação do Operador}

\subsection{Dados Gerais}

\begin{itemize}

    \item Nome do Operador: Vicenzo Soares Minossi
    \item Instituição: PROCERGS - Centro de Tecnologia da Informação e Comunicação do Estado do Rio Grande do Sul S.A
    \item Setor: Divisão Redes e Segurança/Engenharia de Rede
    \item Projeto Associado: RAVEN (Rede Adaptativa de Variáveis em Estado Não-linear)]
    \item Item de Projeto: UBR (Unidade Básica RAVEN)
    \item Variante: "TOTEM"
    \item Versão do Projeto: 1.0
    \item Competência: 2026
    \item Ambiente Operacional: Virtual/Físico
    \item Sistema Operacional: Windows/Linux

\end{itemize}

\vspace{0.5cm}

% =========================================================
% RESUMO / ABSTRACT
% =========================================================

\begin{tcolorbox}[
    colback=white,
    colframe=black,
    width=\textwidth,
    arc=2mm,
    boxrule=0.5pt
]

{\small
\setstretch{1.1}

\section*{\centering Resumo}

\noindent
\textit{
Memorial técnico-científico de dispositivo que interpreta dinamicamente o comportamento de redes de computadores através de microestruturas físicas que formam macroestruturas lógicas discretizadas cujo comportamento pode ser analisado por modelos de grafos; uma arquitetura computacional física, adaptativa e inspirada em redes neurais capaz de executar funções clássicas de gerenciamento e decisão de tráfego em rede.
}

\vspace{0.7cm}

\section*{\centering Abstract}

\noindent
\textit{
Technical-scientific memorial of a device capable of dynamically interpreting the behavior of computer networks through physical microstructures that form discretized logical macrostructures whose behavior can be analyzed through graph models; a physical, adaptive, and neural-network-inspired computational architecture capable of executing classical traffic management and network routing decision functions.
}

}

\end{tcolorbox}

\vspace{1cm}
\newpage
% =========================================================
% QUADRO TEÓRICO
% =========================================================

\section{Quadro Teórico}

\subsection{Fundamentação}

Descrever:

\begin{itemize}

    \item Teoria dos Grafos
    \item Teoria Eletromagnética
    \item Sistemas Dinâmicos
    \item Teoria de Circuitos
    \item Redes de Computadores
    \item Redes Neurais
    \item Escala de modelos de aprendizados de máquina (deep learning e machine learning)
    \item Métodos Computacionais da Física: Cálculo Numérico, Análise de Estabilidade de Método




\end{itemize}

\vspace{0.5cm}

\subsection{Metodologia}

Descrever:
\begin{itemize}

    \item Descrição do comportamento clássico de roteadores em uma rede [DPT]
    \item Estatística de tráfego [DPT]
    \item Métricas de funcionamento[DPT]
    \item Conclusão
    \item Hipótese [DPT]%um algorítimo físico é capaz de gerenciar corretamente uma rede?
    \item Modelo de rede neural proposto [DPT]
    \item Métricas de funcionamento [DPT]
    \item Correlação de métricas: Modelo Clássico vs Modelo Adaptativo
    \item Análise de estabilidade [DPT]
    \item Conclusão [DPT]
    \item Modelo de hardware [DPER]
    \item Aplicações práticas [DPER]

\end{itemize}

% =========================================================
% PSEUDOCÓDIGO
% =========================================================
\section{Introdução}
\subsection{Leonhard Euler: As Sete Pontes de Königsberg}


\section{O Comportamento Clássico de una Rede}
\subsection{A era }
\subsection{Introdução à Redes de Computadores}


\section{Funções Clássicas Fundamentais}
\subsection{Encaminhamento}
\subsection{Descoberta de Topologia}
\subsection{Controle de Fluxo}
\subsection{Segurança}
\subsection{Redundância e Adaptação}\
\subsection{Controle e Telemetria}

\section{Teoria das Redes}
\subsection{A Teoria dos Grafos}
\subsection{Descrição de uma Rede Clássica}
\subsection{Métricas de funcionamento}

\section{Sistemas Dinâmicos}
\subsection{Sistemas não-lineares}
\subsection{Emergência de Computação}






\subsection{Estrutura Lógica}

\begin{lstlisting}[style=ravenstyle, language=Python]


\end{lstlisting}

% =========================================================
% CÓDIGO-FONTE
% =========================================================

\section{Código-Fonte}

\subsection{Inicialização do Sistema}

\begin{lstlisting}[style=ravenstyle, language=Python]



\end{lstlisting}

\vspace{0.5cm}

\subsection{Renderização do Grafo}

\begin{lstlisting}[style=ravenstyle, language=Python]



\end{lstlisting}

% =========================================================
% IDE E BIBLIOTECAS
% =========================================================

\section{Ambiente de Desenvolvimento}

\subsection{IDE Utilizada}

\begin{itemize}

    \item Visual Studio Code
    \item PyCharm
    \item Jupyter Notebook

\end{itemize}

\vspace{0.5cm}

\subsection{Bibliotecas e Módulos}

\begin{itemize}

    \item Python
    \item NetworkX
    \item Matplotlib
    \item Tkinter
    \item PyVis
    \item JSON
    \item NumPy

\end{itemize}

\vspace{0.5cm}

\subsection{Créditos das Ferramentas}

\begin{itemize}

    \item Python Software Foundation
    \item NetworkX Developers
    \item Matplotlib Development Team
    \item PyVis Developers
    \item Open Source Community

\end{itemize}

% =========================================================
% TESTES E RESULTADOS
% =========================================================

\section{Resultados Experimentais}

\subsection{Comportamento Observado}

Inserir:
\begin{itemize}

    \item Logs
    \item Capturas de tela
    \item Métricas
    \item Grafos gerados
    \item Comportamento emergente

\end{itemize}

\vspace{0.5cm}

\subsection{Problemas Encontrados}

Descrever:
\begin{itemize}

    \item Gargalos
    \item Bugs
    \item Saturações
    \item Limitações arquiteturais

\end{itemize}

% =========================================================
% CONCLUSÃO
% =========================================================

\section{Conclusão}

Descrever:
\begin{itemize}

    \item Eficiência do sistema
    \item Aplicabilidade prática
    \item Melhorias futuras
    \item Escalabilidade
    \item Potencial de produção

\end{itemize}

% =========================================================
% BIBLIOGRAFIA
% =========================================================

\section{Bibliografia}

\begin{thebibliography}{99}

\bibitem{graph}
Diestel, Reinhard.
\textit{Graph Theory}.
Springer.

\bibitem{networks}
Barabási, Albert-László.
\textit{Network Science}.
Cambridge University Press.

\bibitem{cyber}
Norbert Wiener.
\textit{Cybernetics}.
MIT Press.

\bibitem{python}
Python Software Foundation.
\textit{Python Documentation}.

\bibitem{networkx}
NetworkX Developers.
\textit{NetworkX Documentation}.

\end{thebibliography}

% =========================================================
% FIM
% =========================================================

\end{document}